
\chapter{Экспериментальное исследование характеристик модели вентилятора}\label{ch:ch3}


	\textit{Третья глава содержит экспериментальное обоснование решения задачи, описание методов экспериментальных исследований, оценку точности, анализ сходимости опытных и теоретических результатов.
Функция третьей, «экспериментальной» главы – конкретизировать обобщенное теоретическое решение задачи, предоставить коэффициенты, полученные опытным путем, привести экспериментальные данные, проверяющие теорию.
Здесь же можно дать описание новых устройств, проверки их работоспособности, новых методов или новой технологии проведения экспериментальных исследований.
Объем третьей главы 30-35 страниц.}


\section{Введение}

Результаты расчётов по разработанной модели течения хорошо согласуются с результатами полученными методом решения системы уравнений RANS методом контрольных объёмов.
Но, несмотря на высокий уровень развития вычислительных методов, решающим фактором применимости моделей является экспериментальная проверка результатов.

Для этого была спроектирована, изготовлена и испытана на аэродинамическом стенде вентилятора с модель вентилятора с меридиональным ускорением потока \autocite{furashov2026}.
Были получены зависимости давления и КПД вентилятора \(\eta\) от производительности \(Q\), а так же параметры по высоте проточной части за рабочим колесом (РК), такие как давление \(p\), величина \(c_2\) и направление \(\alpha_2\) скорости потока.

Сравнение характеристик, полученных методами RANS и экспериментально, позволяет сделать заключение о применимости новой методики проектирования.

\section{Экспериментальная модель}
Модель вентилятора (рисунок~\ref{model}) выполнена методом экструзионной печати (FMD) из полилактида (PLA) и имеет следующие параметры:
\begin{table}[htb]
	\centering
	\caption{Параметры экспериментальной модели}
	\label{tab:ParamExModel}
	\begin{tabular}{l|c|c}
		Диаметр рабочего колеса, м & \(D_\text{К}\) & 0,7\\
		\hline
		Относительный диаметр втулки перед РК & \(\bar{d}_1\) & 0,4 \\
		\hline
		Относительный диаметр втулки между РК и СА & \(\bar{d}_2\) & 0,52 \\
		\hline
		Относительный диаметр втулки за СА & \(\bar{d}_3\) & 0,6 \\
		\hline
		Радиальный зазор лопатки РК, мм & \(s\) & 1 \\
		\hline
		Количество лопаток РК & \(Z_\text{РК}\) & 12 \\
		\hline
		Количество лопаток СА & \(Z_\text{СА}\) & 13 
\end{tabular}
\end{table}

\begin{figure}[htb]
	\centering
	\includegraphics[width = 5cm]{ch3/modelFront}
	\includegraphics[width = 5cm]{ch3/modelBack}
	\caption{Модель вентилятора.}
	\label{model}
\end{figure}

Меридиональное ускорение потока распределено поровну между РК и СА \(c_{3m}/c_{2m}=c_{2m}/c_{1m} \approx 1,14\).

Расчётный коэффициент теоретического давления \(\psi_\text{т} = 1,05\), коэффициент производительности \(\varphi=0,6\).

Профилирование велось на 9 поверхностях тока, делящих проточную часть на струи тока с равным расходом.
Закон профилирования в межвенцовых зазорах "--- постоянная циркуляция (\(\Gamma, \bar{H}_\text{т} = \const\)).
Параметры густоты и углов атак для РК и СА (таблица \ref{foilParam}) задавались для трёх сечений, соответствующих втулке, средней линии вентилятора, периферии.
Промежуточные значения получены интерполяцией.
\begin{table}[htb]
	\caption{Параметры решёток профилей}
	\label{foilParam}
	\begin{center}
		\begin{tabular}{r|c|c|c|c|c|c}

			& \multicolumn{3}{c|}{РК} & \multicolumn{3}{c}{СА} \\
			\hline
			\textnumero~сечения & \(i_\text{рк}\), \textdegree & \(\tau_\text{рк}\) & \(k_\text{рк}\) & \(i_\text{са}\), \textdegree & \(\tau_\text{са}\) & \(k_\text{са}\) \\
			\hline\hline
			(периферия) 1 & \(−2,00\) & 0,90 & 1,10 & 0 & 0,80 & 1,15 \\
			\hline
			2 & \(−1,46\) & 0,96 & 1,09 & 0 & 0,86 & 1,12 \\
			\hline
			3 & \(−0,89\) & 1,03 & 1,08 & 0 & 0,93 & 1,09 \\
			\hline
			4 & \(−0,28\) & 1,10 & 1,07 & 0 & 1,01 & 1,08 \\
			\hline
			(средняя линия) 5 & 0,38 & 1,18 & 1,06 & 0 & 1,09 & 1,07 \\
			\hline
			6 & 1,11 & 1,26 & 1,06 & 0 & 1,18 & 1,09 \\
			\hline
			7 & 1,92 & 1,36 & 1,05 & 0 & 1,27 & 1,12 \\
			\hline
			8 & 2,85 & 1,47 & 1,05 & 0 & 1,37 & 1,17 \\
			\hline
			(втулка) 9 & 4,00 & 1,60 & 1,05 & 0 & 1,49 & 1,26
		\end{tabular}
	\end{center}
\end{table}


\section{Анализ влияния геометрических отклонений лопатки на характеристику}
Для анализа влияния погрешностей формы лопатки при изготовления на аэродинамические характеристики, было проведено контрольное измерение геометрии пера лопатки рабочего колеса методом оптической лазерной триангуляции с точностью до 0,05 мм (рисунок \ref{fig:LTr}).
\begin{figure}[htb]
	\centering
	%\includegraphics{ch3/lop1}
	\includegraphics{ch3/lop2}
	\caption{Отклонение формы лопатки от расчётной (зелёное поле "--- менее 0,1~мм)}
	\label{fig:LTr}
\end{figure}

Изменение коэффициента теоретического напора \(\bar{H}_\text{т}\) было определенно в результате моделирования методом дискретных вихрей (глава \ref{sec:newMDV}) обтекания решётки составленной из профилей с учётом отклонений.

Проведённый анализ показал, что отклонение геометрических параметров приводит к незначительному изменению \(\bar{H}_\text{т}\) не превышающему 0,35\% от номинального значения. 
Углы установки лопаток не отличаются от номинальных более чем на 20', что в среднем не влияет на измерения, а в худшем случае приводящее к ошибке в 0,77\%.
Значение КПД при номинальном режиме работы, как правило, близко к экстремуму и при таких небольших отклонениях \(\bar{H}_\text{т}\), можно считать, остаётся постоянным.
Тогда значение коэффициента напора \(\bar{H}\) изменится пропорционально изменению  \(\bar{H}_\text{т}\), то есть не более 0,35\% от своего номинального значения. 

Влияние шероховатости поверхности на аэродинамические характеристики можно установить по диаграмме на рисунке \ref{fig:Prandl1934}, на котором отображено влияние высоты шероховатости \(k_s\) на коэффициент сопротивления пластины \(c_f\) длиной \(l\) при турбулентном течении от числа Рейнольдса \(U_\infty l/\nu\), где \(\nu)\) "--- кинематическая вязкость, м\textsuperscript{2}/с.
\begin{figure}[htb]
	\centering
	\includegraphics[width=\textwidth]{ch3/Prandl1934}
	\caption{Влияние шероховатости на коэффициент сопротивления пластины при турбулентном течении \autocite{Prandl1934, Schlichting1974}.}
	\label{fig:Prandl1934}
\end{figure}

При расчётном значении числа Рейнольдса в рабочем колесе \(\Re_{w_1}=1,88\times10^5\) и средней длине хорды профиля \(b_\text{ср}=145\)~мм гидравлическая гладкость поверхности пластины достигается при высоте шероховатости ниже 0,15~мм.
Шероховатость при экструзионной печати не анизотропна, и максимальная высота  достигается в направлении поперёк движения экструдера.
При толщине слоя печати 0,25~мм, в крайнем случае, высота шероховатости составляет лишь 0,125~мм, меньше допустимой.

Таким образом, в рамках поставленных задач, влияние погрешностей изготовления на итоговые экспериментальные данные являются пренебрежимо малым.

\section{Экспериментальный стенд и методика измерений}
Исследования проводились в соответствии с ГОСТ 10921 — 2017 <<Вентиляторы радиальные и осевые. Методы аэродинамических испытаний>> на стенде К-1 в Московском комплексе Центрального аэрогидродинамического института имени профессора Н.\, Е.~Жуковского (МК~ЦАГИ).

Стенд (рисунок \ref{stend}) представляет собой камеру всасывания 1, замеряемое разряжение в которой создаёт действующая модель вентилятора 2.
В камеру всасывания воздух поступает из помещения лаборатории через расходомерный коллектор 3 и выравнивается специальными сетками и хонейкомбом 4.
Электродвигатель 5 приводит во вращение ротор модели вентилятора.
Момент на валу и частота вращения измеряется датчиком 6.
Расход воздуха регулируется дросселем 7 и вспомогательным вентилятором наддува 8.
\begin{figure}[htb]
	\centering
	\includegraphics[scale=1.17]{ch3/stand}
	\caption{Схема экспериментального стенда.}
	\label{stend}
\end{figure}


\subsection{Интегральные параметры работы вентилятора}

Стенд позволяет измерять следующие параметры потока: давление \(p_\text{м}\) в камере всасывания, разряжение в расходомерном коллекторе \(p_c\), частоту вращения модели вентилятора \(n\), момент на валу модели \(M\).

Параметры атмосферного воздуха определены с погрешностью для относительной влажности \(\varphi_a \pm 2\)\%; температуры \(T_a \pm 0,2\)~K; давления \(p_a \pm 300\)~Па. 

Давление измеряется наклонными спиртовыми манометрами. 
При различных положениях синус угла между мерной шкалой и горизонтом составляет 1, 0,5, 0,25 и 0,125.
Цена деления шкалы манометра 1~мм, диапазон шкалы до 220~мм.
Измеритель крутящего момента имеет диапазон измерений \(\pm50\)~Н\(\cdot\)м и класс точности~0,3.

Максимальное полное давление вентилятора \(p_\nu\) в эксперименте составило 145~Па, число Маха, 0,07. 
Так как \(p_\nu<3\)~кПа и число Маха меньше 0,15 "--- влиянием сжимаемости воздуха на характеристики было пренебрежено.

Плотность атмосферного воздуха определяется по уравнению состояния идеального газа
\begin{equation}
	\rho_a= \frac{p_a}{R_a t_a},
	\label{eq:startExp}
\end{equation}
\begin{eqexpl}
\item{\(p_a\)} атмосферное давление, Па;
\item{\(R_a\)}[=]\(287\)~Дж/кг~K газовая постоянная воздуха;
\item{\(t_a\)}  температура атмосферного воздуха, К.
\end{eqexpl}

Производительность вентилятора определяется по разряжению в расходомером коллекторе \(p_c\), Па
\begin{equation}
	Q = \alpha_\text{к} \varepsilon \frac{\pi d^2}{4} \sqrt{\frac{2p_c}{\rho_a}},
\end{equation}
\begin{eqexpl}
\item{\(\alpha_\text{к}\)} коэффициент расхода расходомерного коллектора;
\item{\(\varepsilon\)} коэффициент расширения воздуха в коллекторе;
\item{\(d\)} диаметр проточной части коллектора в мерном сечении, м.
\end{eqexpl}

Коэффициент \(\varepsilon\) вычисляют по формуле
\begin{equation}
	\varepsilon=1-\frac{3}{2\kappa} \frac{p_c}{p_a},
\end{equation}
\begin{eqexpl}[12 mm]
	\item{\(\kappa=1,4\)} показатель адиабаты воздуха.
\end{eqexpl}


Динамическое давление вентилятора подсчитывается по скорости потока на выходе из проточной части, полученной по производительности вентилятора отнесенной к площади проточной части.
\begin{equation}
	p_{d\nu}=\frac{\rho_a}{2} \left(\frac{Q}{F_2}\right)^2,
\end{equation}
\begin{eqexpl}
	\item{\(F_2\)} площадь проточной части на выходе из вентилятора, \(\text{м}^2\).
\end{eqexpl}


Статическое давление вентилятора определяют по формуле
\begin{eqnarray}
	p_{s\nu}=p_\text{м}-p_{d_1},
\end{eqnarray}
\begin{eqexpl}
	\item{\(p_{d_1}\)} динамическое давление потока перед вентилятором без учёта сжимаемости, Па,
	\begin{equation*}
		p_{d_1} = \frac{\rho_a}{2} \left(\frac{4Q}{\pi D_1^2} \right)^2,
	\end{equation*}
	\item{\(D_1\)} диаметр успокоительной камеры, м.
\end{eqexpl}

Полное давление вентилятора \(p_{\nu}\) равно сумме статического и полного давления вентилятора
\begin{equation}
	p_\nu=p_{d\nu}+p_{s\nu}.
\end{equation}

Потребляемая вентилятором мощность \(N\), Вт, определялась с помощью измерителя крутящего момента и частоты вращения рабочего колеса \(n\), об/мин
\begin{equation}
	N=M \frac{\pi n}{30}.
\end{equation}

Полезная мощность по полному давлению определялась как
\begin{equation}
	N_\nu = p_\nu Q.
	\label{eq:stopExp}
\end{equation}

Характеристики вентилятора были пересчитаны в безразмерных величинах КПД \(\eta\) и коэффициента теоретического напора \(\bar{H}_\text{т}\) от коэффициента расхода на выходе из спрямляющего аппарата \(\bar{c}_{3m}\) (раздел \ref{sec:ch1/sec1}).

\subsection{Анализ неопределённости измерений}
Анализ непределённости результатов измерений основан на допущениях о линейной зависимости значении исследуемых величин вида \(Y(x_1, x_2, \dots x_n)\) от отклонений параметров \(\Delta_{x_1}, \Delta_{x_2}, \dots \Delta_{x_n}\):
\begin{equation}
	\sigma_Y(x_1, x_2, \dots x_n) = \sqrt{\left(\frac{\partial Y}{\partial{x_1}}\sigma_{x_1}\right)^2 + \left(\frac{\partial Y}{\partial{x_2}}\sigma_{x_2}\right)^2 + \dots + \left(\frac{\partial Y}{\partial{x_n}}\sigma_{x_n}\right)^2},
\end{equation}
\begin{eqexpl}
	\item{\(\sigma_Y\)} среднее квадратическое отклонение величины \(Y\);
	\item{\(\sigma_x\)} среднее квадратическое отклонение параметра \(x\).
\end{eqexpl}

При анализе относительных средних квадратичных величин \(\epsilon_Y = \sigma_Y/Y\) в зависимости от независимых переменных \(x_{1,2}\) справедливы следующие выражения:
\begin{itemize}
\item \(n\epsilon_x\) "--- для величины \(x^n\);
\item \(\sqrt{\epsilon_{x_1}^2+\epsilon_{x_2}^2}\) "--- для величины \(x_1 x_2\);
\item \(\sqrt{\frac{\epsilon_{x_1}^2 {x_1}^2 + \epsilon_{x_2}^2 {x_2}^2}{x_1^2 \pm x_2^2}}\) "--- для величины \(x_1 \pm x_2\).
\end{itemize}
Анализ зависимостей \ref{eq:startExp}--\ref{eq:stopExp} \autocite{Kersten1993} позволил определить неопределённость измерений параметров в виде относительных средний квадратичных отклонений (ОСКО) следующим образом:

ОСКО расхода воздуха:
\begin{equation}
	\epsilon_Q = \sqrt{\epsilon_\alpha^2 + 4\epsilon_d^2 + 0,25(\epsilon_{\rho_a}^2 + \epsilon_{p_\text{c}}^2)},
\end{equation}
\begin{eqexpl}
	\item{\(\epsilon_\alpha\)} ОСКО коэффициента расходомерного коллектора \(\epsilon_\alpha = 0,5\%\);
	\item{\(\epsilon_d\)} ОСКО диаметра расходомерного коллектора \(\epsilon_d = 0,1\%\);
	\item{\(\epsilon_{\rho_a}\)} ОСКО плотности атмосферного воздуха;
	\item{\(\epsilon_{p_\text{c}}\)} ОСКО давления в расходомерном коллекторе.
\end{eqexpl}

ОСКО динамического давления вентилятора:
\begin{equation}
	\epsilon_{p_{d\nu}} = \sqrt{\epsilon_\Delta^2 + 4\epsilon_Q^2 + 16\epsilon_D^2},
\end{equation}
\begin{eqexpl}
	\item{\(\epsilon_d\)} ОСКО диаметра корпуса вентилятора \(\epsilon_d = 0,1\%\).
\end{eqexpl}

ОСКО статического давления вентилятора:
\begin{equation}
	\epsilon_{p_{s\nu}} = \sqrt{\epsilon_{p_\text{м}}^2 \frac{p_\text{м}}{p_{s\nu}}},
\end{equation}
\begin{eqexpl}
	\item{\(\epsilon_{p_\text{м}}\)} ОСКО разности статического давления за вентилятором и перед ним;
\end{eqexpl}

ОСКО полного давления вентилятора:
\begin{equation}
	\epsilon_{p_{\nu}} = \sqrt{\epsilon_{p_{s\nu}}^2 \frac{p_{s\nu}}{p_{\nu}}+\epsilon_{p_{d\nu}}^2 \frac{p_{d\nu}}{p_{\nu}}},
\end{equation}

ОСКО мощности:
\begin{equation}
	\epsilon_{N} = \sqrt{\epsilon_{M}^2 +\epsilon_{n}^2},
\end{equation}
\begin{eqexpl}
	\item{\(\epsilon_{M}\)} ОСКО момента на валу;
	\item{\(\epsilon_{n}\)} ОСКО частоты вращения вала на валу;
\end{eqexpl}

ОСКО КПД:
\begin{equation}
	\epsilon_{\eta} = \sqrt{\epsilon_{p_\nu}^2 +\epsilon_{Q}^2 + \epsilon_{N}^2}.
\end{equation}

\subsection{Параметры потока по высоте проточной части}
Исследование параметров потока по высоте проточной части за рабочим колесом проводилось ориентируемым трёхканальным насадком (рисунок \ref{nasadok}) \autocite{Peschechonov1962a, Petunin1972}.
Насадок крепится на координатном устройстве, задающем его радиальное и угловое положение.
Боковые приёмники давления (1 и 2) позволяют установить насадок по потоку, и вместе с центральный приёмником (0) "--- определить полное и статическое давление.
Загромождение насадком в эксперимента не превышало \(0,36\%\) от площади проточной части.
\begin{figure}[htb]
	\centering
	\includegraphics{ch3/nasadok}
	\caption{Схема трёхканального насадка.}
	\label{nasadok}
\end{figure}

Скорость потока \(c\), м/с, определяется по разнице давлений между центральным и боковыми приёмниками давлений \(\Delta p_{01}\), Па,
\begin{equation}
	c=\sqrt{2\zeta \frac{\Delta p_{01}}{\rho}},
\end{equation}
\begin{eqexpl}[7 mm]
	\item{\(\zeta\)}[=]\(0,734\) "--- коэффициент насадка, определённый при тарировке в аэродинамической трубе Т-23 в МК ЦАГИ \autocite{Zamolodchikov2019}.
\end{eqexpl}

Коэффициент измеренной меридиональной компоненты скорости определяется по значению скорости и углу \(\alpha\) между направлением потока и меридиональной плоскостью
\begin{equation}
	\bar{c}_{m \, \text{изм}} =\frac{c \cos{\alpha}}{U}.
\end{equation}

Коэффициент измеренной окружной скорости
\begin{equation}
	\bar{c}_{u \, \text{изм}}=\frac{c \sin{\alpha}}{U}.
\end{equation}

Коэффициент напора
\begin{equation}
	\bar{H}=\frac{p^{*}}{\rho U^2}.
\end{equation}

Измеренные насадком значения \(\bar{c}_{m \,\text{изм}}\) и \(\bar{c}_{u \,\text{изм}}\) уточняются по интегральным значениям коэффициентов производительности и теоретического напора

Так как через насадок проходят возмущения от лопаток, измеренные значения менее надёжны чем интегральные параметры \(\bar{H}_\text{т}\) и \(\varphi\), полученные на стенде. 
Полученные распределения \(\bar{c}_{m}(\bar{r})\) и \(\bar{c}_{u}(\bar{r})\) нормируются с учётом значений \(\bar{H}_\text{т}\) и \(\varphi\):
\begin{gather}
	\bar{c}_m = k_m \bar{c}_{m \, \text{изм}}, \\
	\bar{c}_u = k_u \bar{c}_{u \, \text{изм}},
\end{gather}
\begin{eqexpl}
	\item{\(k_m\)}[=]\(\varphi/(2 \int_{\bar{d}}^{1} \bar{c}_{m \, \text{изм}} \bar{r} \, d\bar{r})\) "--- коэффициент осреднения меридиональной компоненты скорости,

	\item{\(k_u\)}[=]\(\varphi \bar{H}_\text{т}/(2 \int_{\bar{d}}^{1} \bar{c}_{u \, \text{изм}} \bar{c}_m \bar{r}^2 d\bar{r})\) "---  коэффициент осреднения окружной компоненты скорости,

	\item{\(\bar{r}\)}[=]\(r/R_\text{к}\) "--- относительный радиус струйки тока.
\end{eqexpl}

Коэффициент теоретического напора определяется по разнице циркуляций за (2) и перед (1) РК
\begin{equation}
	\bar{H}_\text{т} = (\bar{c}_{2u} \, \bar{r}_2 - \bar{c}_{1u} \, \bar{r}_1),
\end{equation}
\begin{eqexpl}
	\item{\(\bar{r}\)}[=]\(r/R_\text{к}\) "--- относительный радиус струйки тока.
 \end{eqexpl}

\section{Результаты эксперимента}
\subsection{Характеристика вентилятора}
Аэродинамическая характеристика модели вентилятора (рисунок~\ref{charact}) была получены при частоте вращения \(n = 450\)~об/мин с окружной скоростью периферии рабочего колеса \(U \approx 16,5\)~м/с.
Число Рейнольдса на номинальном режиме работы, посчитанное по относительной скорости перед решёткой \ref{eq:Re_w}, составило для РК \(\Re_\text{к} = 1,88 \times 10^5\), а для СА определенное по абсолютной скорости \(\Re_\text{са} = 1,73 \times 10^5\).
\begin{figure}[htb]
	\centering
	\includegraphics{ch3/characteristic}
	\caption{Аэродинамическая характеристика вентилятора.}
	\label{charact}
\end{figure}

Характеристика состоит из двух ветвей "--- правой и левой "--- связанных между собой петлёй гестерезиса. 
Правая является рабочей и обеспечивает высокий КПД вблизи расчётного режима работы. 
Нерабочая левая ветвь имеет низкую эффективность из-за развитых отрывных явлений.
Переход между ветвями сопровождается скачкообразным изменением параметров от \(\bar{c}_{3m} = 0,63\) до 0,75 слева направо.
С правой на левую скачком от 0,63 до 0,42.

При расчётном коэффициенте производительности \(\bar{c}_{3m}=0,94\), коэффициент теоретического давления составил \(\bar{H}_\text{т} = 0,531 \pm 0,006\), КПД \(\eta = 0,841 \pm 0,011\).


%	Такое значение КПД относительно высоко для осевых вентиляторов при данных числах Рейнольдса.
%	Для эффективного осевого вентилятора с цилиндрической проточной частью и максимальным значением \(\eta=0,87\), при данном числе Рейнольдса \(\eta\) не превысит 0,79 \autocite{Brusilovskiy1984}.


\subsection{Влияние числа Рейнольдса на величину потерь полного давления}
Характеристика исследована при числах Рейнольдса~\(\Re_{w_1}\) значительно меньших критического значения автомодельности для вентиляторов с меридиональным ускорением (\(\Re_\text{а} = 9\times10^5\))~\autocite{Brusilovskiy1962}.

Изменение величины потерь полного давления \(1–\eta\), по аналогии с течением в каналах, описывается степенной зависимостью:
\begin{equation}
	\left(\frac{\Re_1}{\Re_2}\right)^k=\frac{1-\eta_1}{1-\eta_2}, \label{eq:ReEta}
\end{equation}
\begin{eqexpl}
	\item{\(k\)}[=]\(-0.2\) "--- показатель степени для турбулентного течения, 
	\item{\(k\)}[=]\(-0.5\) "--- для ламинарного. 
\end{eqexpl}
Однако в реальности взаимодействие пограничных слоёв и вторичных течений в проточной части приводит к искажению этой зависимости. 
От числа Re зависят в первую очередь профильные потери. 
Вторичные потери, образованные поворотом потока в межлопаточных каналах, напротив, от Re зависят мало \autocite{Dovjik1968}. 
По этой причине показатель степени k плавно меняется с увеличением  не достигая \(-0.2\) \autocite{Brusilovskiy1984}. 

При расчёте методом RANS было принято допущение о полностью турбулентном пограничном слое. 
При \(\Re_{w_1}<\Re_\text{а}\) влияние на результат этого допущения может стать существенным и непредсказуемым, так как затягивание ламинарно-турбулентного перехода пограничного слоя приводит к двум противодействующим эффектам.
С одной стороны, к снижению потерь полного давления за счёт увеличения длины ламинарной зоны. 
С другой, напротив, привести к росту потерь из-за большей склонности к отрыву ламинарного пограничного слоя, по сравнению с турбулентным. 

Для проверки потерь полного давления определяемых моделированием потока методом RANS проведено численное и экспериментальное исследование изменения потерь \(1-\eta\) от чисел Рейнольдса~\(\Re_{w_1}\) на номинальном режиме работы.
Изменение \(\Re_{w_1}\) в диапазоне от \(1\times10^5\) до \(4\times10^5\) было получено варьированием частоты вращения рабочего колеса \(n\).

Результаты приведены на рисунке \ref{fig:re}. 
Столбиками обозначены пределы неопределённости результатов эксперимента выраженные средним квадратичным отклонением величины \(\eta\).
\begin{figure}[htb]
	\centering
	\includegraphics{ch3/Re}
	\caption{Сравнение относительных потерь полного давления (\(1-\eta\)) модели вентилятора определённых экспериментально и численно (RANS) в диапазоне чисел \(\Re_\text{РК}\) от \(1\times10^5\) до \(4\times10^5\).}
	\label{fig:re}
\end{figure}

Методом RANS достаточно точно предсказано изменение \(1-\eta\).
Все полученные результаты находятся в пределах доверительного интервала эксперимента.
Значения \(\eta\) можно аппроксимировать зависимостью \ref{eq:ReEta} с коэффициентом \(k = -0,289\).
Эта величина, в исследованном диапазоне чисел Re, соответствует данным \autocite{Brusilovskiy1984} для вентилятора РК+СА.

\subsection{Параметры по высоте проточной части за РК}
По высоте проточной части \(h\) в сечении за РК CFD-моделированием и экспериментально (цилиндрическим насадком) определены параметры потока \(\bar{c}_{2m}(\bar{h})\), \(\bar{H}(\bar{h})\) и \(\bar{H}_\text{т}(\bar{h})\), где \(\bar{h} = h/(1-d_2)\) "--- относительная высота проточной части.
Исследовано четырё режима работы, при \(\bar{c}_{3m}\) равном 0,94 (номинальный), 1,25 и 0,63 на рабочей и нерабочей ветви.

На номинальном режиме работы (\(\bar{c}_{3m} = 0,94\)) параметры в средней части проточной части (\(\bar{h}\) от 0,2 до 0,8) близки к расчётным (рисунок \ref{fig:c2m_exp} и \ref{fig:Ht_exp}).
При увеличенной производительности (\(\bar{c}_{3m} = 1,25\)) меридиональная компонента скорости остаётся примерно постоянной по высоте проточной части, а \(\bar{H}\) и \(\bar{H}_\text{т}\) уменьшается к периферии.
\begin{figure}[htb]
	\centering
	\includegraphics[page=1]{ch3/c2m_exp}
	\caption{Распределение экспериментального \(\bar{c}_{2m}\) по относительной высоте проточной части \(\bar{h}\).}
	\label{fig:c2m_exp}
\end{figure}
\begin{figure}[htb]
	\centering
	\includegraphics[page=2]{ch3/c2m_exp}
	\caption{Распределения экспериментальных \(\bar{H}_\text{т}\) и \(\bar{H}\) по относительной высоте проточной части \(\bar{h}\).}
	\label{fig:Ht_exp}
\end{figure}

При уменьшении производительности к границе перехода с правой на левую ветвь характеристики (\(\bar{c}_{3m} = 0,63\)) возрастает \(\bar{H}_\text{т}\) у периферии, а \(\bar{H}\) снижаются.
Переход на левую ветвь сопровождается снижением \(\bar{H}\) по всей высоте и снижением \(\bar{H}_\text{т}\) означающем отрыв потока, особенно резким у периферии.
Распределения \(\bar{c}_{2m}\) практически не изменяются.

\section{Сравнение с расчётной характеристикой}
Моделированием были получены значения только на рабочей ветви.

Относительные разницы численного моделирования интегральных величин \(\eta\) и \(\bar{H}_\text{т}\) подсчитывались по формуле:
\begin{equation}
	\bar{\Delta} = 1 - \frac{x_\text{CFD}}{x_\text{эксп}},
\end{equation}
\begin{eqexpl}[12 mm]
	\item{\(x_\text{CFD}\)} значение полученное численным методом; 
	\item{\(x_\text{эксп}\)} экспериментальное.
\end{eqexpl}

Значения \(\bar{\Delta}\) и величины относительной неопределённости экспериментальных значений \(\varepsilon_\text{эксп}\) в зависимости от коэффициента расхода \(\bar{c}_{3m}\) приведены на рисунке \ref{fig:errCFX}.
\begin{figure}[htb]
	\centering
	\includegraphics{ch3/err}
	\caption{Область расширенной неопределённости эксперимента \(\epsilon\) и относительная разница \(\bar{\Delta}\) экспериментальных и расчётных характеристик, полученных квазитрёхмерным методом (Q3D) и RANS.}
	\label{fig:errCFX}
\end{figure}

Значения \(\bar{H}_\text{т}\), полученные при RANS-моделировании рабочей ветви характеристики, оказались несколько заниженными относительно экспериментальных, однако остались в пределах расширенной неопределённости эксперимента \(P=0,95\)~(\(\pm2\varepsilon\)).
Значения \(\eta\) совпали с экспериментальными в пределах \(0,5\%\) при \(\varphi\) от 0,62 до 1,1, но при 1,25 показали значительно завышенный результат (на \(2,7\%\)).
%\begin{figure}[htb]
%	\centering
%	\includegraphics[width=10cm]{ch3/characteristic_compare}
%	\caption{Сравнение экспериментальных и расчётных характеристик}
%	\label{fig:characteristicCFX}
%\end{figure}

При CFD-моделировании (\(\bar{c}_{3m} = 1,25\)) не наблюдается столь существенного отставания \(\bar{H}\) от \(\bar{H}_\text{т}\) начиная с середины лопатки, как в эксперименте.
Это приводит к существенной ошибке в определении интегрального КПД.

\begin{figure}[htb]
	\centering
	\includegraphics[page=1]{ch3/c2m_compare}
	\caption{Сравнение расчётных и экспериментальных \(\bar{c}_{2m}\) по относительной высоте проточной части \(\bar{h}\).}
	\label{fig:c2m_compare}
\end{figure}
\begin{figure}[htb]
	\centering
	\includegraphics[page=2]{ch3/c2m_compare}
	\caption{Сравнение расчётных и экспериментальных \(\bar{H}_\text{т}\) по относительной высоте проточной части \(\bar{h}\).}
	\label{fig:c2m}
\end{figure}
Наибольшее отличие экспериментальных распределений от полученных CFD-моделированием в периферийной части.
След от радиального зазора лопатки имеет пространственную форму и вращается с рабочим колесом, экспериментальные значения полученные неподвижным насадком, в данном случае, ненадёжны.

\section{Заключение}
Проведённое исследование аэродинамических характеристик вентилятора с меридиональным ускорением тремя методами позволило получить следующие результаты:
\begin{enumerate}
	\item Модель вентилятора с меридиональными ускорением потока, спроектированного по разработанному методу, была испытана на стенде для определения его интегральных аэродинамических характеристик.
	Коэффициент теоретического напора \(\bar{H}_\text{т}\) на номинальном режиме работы (\(\bar{c}_{3m}=0,94\)) составил \(0,531\). 
	В сравнении с расчётных значением 0,525, относительная разница составляет 0,9\%;

	\item Моделирование потока методом RANS с замыканием моделью турбулентности \(k-\omega\) позволило предсказать характеристику вентилятора с точностью относительно эксперимента по КПД "--- менее 1\% на рабочей ветви при \(\bar{c}_{3m}\) от 0,62 до 1,1; по \(\bar{H}_\text{т}\)  "--- 2,1\%.
	При \(\bar{c}_{3m}\) превышающих 1,1 результаты моделирования значительно расходятся с экспериментальными, ошибка КПД превысила доверительный интервал эксперимента и составила более 2,7\%;

	\item С помощью моделирования методом RANS было предсказано изменение потерь полного давления вентилятора на номинальном режиме работы при числах Рейнольдса меньше автомодельного.
	Полученная зависимость было подтверждена экспериментально и хорошо аппроксимируется в исследованных диапазонах (\(\Re_{w_1} = 1 \dots 4 \times 10^5\)) степенной зависимостью с показателем степени \(k=-0,289\);

	\item Результаты экспериментального и численного исследования параметров потока по высоте проточной части хорошо совпали на исследованных режимах работы (\(\bar{c}_{3m}=0,63; 0,94; 1,25\)).
	Относительное среднее квадратичное отклонение результатов составило ...
	%В периферийной части эксперимент показал снижение параметров, не наблюдаемое при численном моделировании, однако взаимодействие измерительного насадка и следа от радиального зазора лопаток рабочего колеса снижает достоверность полученных экспериментальных данных.
	%Переход с рабочей ветви характеристики на нерабочую (при \(\varphi=0,4\)) сопровождается резким снижением коэффициента давления \(\psi\) по всей высоте проточной части, и теоретического давления \(\psi_\text{т}\) у периферии, свидетельствующем об отрыве потока.
\end{enumerate}

Разработанный квазитрёхмерный метод моделирования потока в вентилятора и основанная на нём методика профилирования лопаток в вентиляторах с меридиональным ускорением потока подтверждёны экспериментально.

\FloatBarrier
